\section{Related Work}
\paragraph{Temporal summaries and spatial linearization.}
Generalized space--time cubes describe temporal visualizations through transformations of space and time~\cite{Bach:2017:STC}. Studies of animation and trend analysis show how presentation and comparison tasks affect the usefulness of animated and static displays~\cite{Tversky:2002:Animation,Robertson:2008:Animation}. Static temporal maps place distant time steps in one view, with their interpretation determined by the compressed spatial axis. Hovm\"oller diagrams retain a selected spatial coordinate~\cite{Hovmoller:1949}; geographic orderings and context-based space-filling curves construct spatial traversals~\cite{Guo:2006:SpatialOrder,Dafner:2000:CSFC}. Dynamic Volume Lines apply space-filling traversals to volumetric comparison~\cite{Weissenboeck:2019:DVL}, while one-dimensional projections and data-driven curves adapt the ordering to spatiotemporal fields~\cite{Franke:2021:1DP,Zhou:2021:DSF}. MotionRugs instead organize moving point features~\cite{Buchmueller:2019:MR}. RG-MTM gives the axis a fixed projection target and quantifies the displacement introduced by hierarchical feature grouping.

\paragraph{Feature tracking and evolving hierarchies.}
Feature tracking connects spatial structures across time, including evolving regions and turbulent three-dimensional features~\cite{Samtaney:1994:Tracking,Silver:1997:Tracking}. Topology-based approaches organize such correspondences through multiscale structures, global scalar-feature matching, and time-varying merge trees~\cite{Bremer:2010:Burning,Saikia:2017:GFT,Oesterling:2017:TVMT}. Dynamic tracking graphs support exploration of temporal relationships~\cite{Widanagamaachchi:2012:DTG}; nested and dynamic nested tracking graphs add the hierarchical relationships of scalar features~\cite{Lukasczyk:2017:NTG,Lukasczyk:2020:DNTG}. Temporal Treemaps summarize evolving trees as time-aligned bands~\cite{Koepp:2019:TT}, and Fuzzy Contour Trees align multiple trees in a joint layout~\cite{Lohfink:2020:FCT}. These methods establish feature identity and hierarchical context across observations. RG-MTM uses support-overlap matches as input and addresses the subsequent placement of matched extrema against a fixed reference.

\paragraph{Topology-based scalar maps.}
Topological landscapes represent scalar connectivity through terrain, with geometry-preserving extensions introducing proximity between features~\cite{Weber:2007:TL,Beketayev:2012:GPTL}. TMTM traverses augmented merge trees to keep features intact and optimizes subtree overlap between adjacent time steps~\cite{Koepp:2023:TMTM}. ST-MTM separates a hierarchy-constrained skeleton from leaf-arc and path filling, improving geometric distances and temporal coherence~\cite{STMTM:2026}. Section~\ref{sec:prior_pipelines} and Fig.~\ref{fig:prior_work} summarize both workflows before defining our reference-guided extension. RG-MTM reuses this scalar/path filling and separates the original tree, which supplies merge targets, from the ordering tree, which specifies admissible permutations. Pairwise distances can remain unchanged under a global translation or reflection, whereas fixed-reference residuals detect either transformation. RG-MTM therefore adds an absolute positional target to the geometry-aware design space. Its certificate measures the reference-position obstruction, and its scalar correction minimizes merge-level error for the selected order and fixed extremum values.

\paragraph{Ordering, simplification, and tree comparison.}
Optimal leaf ordering searches tree-consistent permutations~\cite{BarJoseph:2001:OLO}. Separation-constrained layout and interval separation control overlap and displacement~\cite{Dwyer:2006:IPSEP,Li:2019:SOI}. RG-MTM couples these constraints with reference coordinates, interval widths, and anchor slack. Persistence simplification removes scalar features~\cite{Edelsbrunner:2002:TPS,Tierny:2012:GTS}; our ordering relaxation retains the leaves and releases selected contiguity constraints. Topological comparison methods span several descriptors and matching assumptions~\cite{Yan:2021:Survey}, including labelled cophenetic/interleaving comparisons~\cite{Munch:2019:LINF,Yan:2022:GAMT}, Wasserstein merge-tree comparisons~\cite{Pont:2022:WMT}, and branch-decomposition-independent edit distances~\cite{Wetzels:2022:BDI}. Our discrepancy uses fixed leaf correspondences to measure the scalar change introduced by an order and its filled profile.

\paragraph{Layout quality.}
Stress and trustworthiness assess distance and neighbourhood preservation~\cite{Smelser:2025:Scale,Venna:2001:NP,Nonato:2019:MP}; ST-MTM also evaluates temporal displacement~\cite{STMTM:2026}. We combine these measures with reference-order error, actual anchor displacement, and rendered merge-level error, distinguishing the spatial, temporal, and topological properties represented by the map.
